% ============================================================
\section{Detalle de módulos}\label{sec:modulos}
% ============================================================
Se describen las características principales de cada módulo y la forma en que se manejan.
El detalle completo está en las hojas de datos citadas.

% ------------------------------------------------------------
\subsection{Controlador: STM32F103C8T6 (\emph{Blue Pill})}
Microcontrolador de 32 bits con núcleo ARM Cortex-M3 de STMicroelectronics~\cite{stm32ds,rm0008}:
\SI{64}{\kilo\byte} de Flash, \SI{20}{\kilo\byte} de SRAM, frecuencia máxima de
\SI{72}{\mega\hertz}, 4 temporizadores de 16 bits, controlador de interrupciones NVIC con
controlador de interrupciones externas EXTI (\emph{External Interrupt}), interfaces
I\textsuperscript{2}C, SPI, USART y un periférico USB 2.0
\emph{Full Speed} nativo. La placa \emph{Blue Pill} incorpora un cristal de \SI{8}{\mega\hertz},
regulador de \SI{3,3}{\volt}, conector micro-USB (conectado a PA11/PA12) y el LED de usuario en
PC13. Se programa y depura por SWD con un ST-Link desde STM32CubeIDE, usando la biblioteca
HAL~\cite{hal}.

Se eligió frente al Arduino UNO del diseño original por su mayor capacidad de cómputo, su
USB nativo (sin conversor USB-serie), la cantidad de temporizadores y la flexibilidad del
EXTI para asignar una interrupción independiente a cada sensor.

\fig[0.6\textwidth]{pinout_stm32}{Asignación de pines del STM32F103C8T6 en STM32CubeMX.}{fig:pinout}{\texttt{Documentación/conf\_pines\_stm32f103c8t6.png}}

% ------------------------------------------------------------
\subsection{Motores paso a paso}
Los tres ejes se accionan con motores paso a paso de tamaño NEMA~17 (\SI{42}{\milli\meter})
recuperados de fotocopiadoras fuera de servicio: dos bipolares y uno unipolar, este último
adaptado para funcionar como bipolar (se dejan sin conectar las tomas centrales de sus
bobinados, de modo que cada fase se excita completa con polaridad reversible). La
Tabla~\ref{tab:motores} resume los datos de las hojas de datos disponibles.

\begin{table}[H]
\centering\small
\caption{Motores paso a paso utilizados.}\label{tab:motores}
\begin{tabularx}{\textwidth}{@{}>{\raggedright\arraybackslash}X >{\raggedright\arraybackslash}p{2.6cm} cccc@{}}
\toprule
\textbf{Modelo} & \textbf{Tipo} & \textbf{Paso} & \textbf{Corriente} & \textbf{Resist.} & \textbf{Torque ret.}\\
\midrule
NMB 17PM-K049U~\cite{nmb17pm} & Híbrido, unipolar (6 hilos) & \ang{1,8} & \SI{1}{\ampere} & \SI{3,1}{\ohm} & \SI{200}{\milli\newton\meter}\\
NMB PM42L-048~\cite{nmbpm42} & Imán permanente, bipolar & \ang{7,5} & \SI{0,6}{\ampere} & \SI{7}{\ohm} & ---\\
\TODO{modelo del 2.º bipolar} & Bipolar & \ang{1,8} & --- & --- & ---\\
\bottomrule
\end{tabularx}
\end{table}

\TODO{Confirmar qué motor acciona cada eje y completar la tabla (pendiente de relevar el modelo
exacto). Datos provisionales anotados en el relevamiento: el motor unipolar (adaptado a bipolar)
estaría en el eje X (base), mientras que los ejes Y y Z llevarían los dos motores bipolares.
Marcas parciales registradas: ``17PM / J344 / G2VS'' y ``17PU / H502 / P5VS''. El firmware trata
X (base) e Y (hombro) como motores de \ang{1,8} (200 pasos/vuelta) y Z (codo) como de \ang{3,75}
(medio paso de un motor de \ang{7,5}, 48 pasos/vuelta); verificar el paso real de Z contra su
hoja de datos.}

\fig[0.7\textwidth]{motores}{Motores paso a paso recuperados utilizados en el proyecto.}{fig:motores}{fotos de \texttt{Documentación/Fotos/motores}}

% ------------------------------------------------------------
\subsection{Driver de motor paso a paso: Pololu A4988}
El A4988 de Allegro~\cite{a4988,pololu} es un driver de micropasos para motores bipolares con
doble puente H de MOSFET y regulación de corriente por \emph{chopper} (PWM de corriente
constante). Admite una tensión de motor de \SIrange{8}{35}{\volt}, hasta \SI{2}{\ampere} por
fase con disipación adecuada (aprox.\ \SI{1}{\ampere} sin disipador) y cinco resoluciones
(paso completo, 1/2, 1/4, 1/8 y 1/16) seleccionadas con MS1--MS3. Incluye protección térmica,
por subtensión y contra corriente cruzada.

Se maneja con dos señales: cada flanco ascendente en \textbf{STEP} avanza un (micro)paso y el
nivel de \textbf{DIR} fija el sentido. La entrada \textbf{ENABLE} es activa en bajo; en \tails{}
las tres entradas están unidas a PB8, de modo que un único pin habilita o libera el torque de
los tres ejes. Los tiempos mínimos que exige el driver (pulso STEP de \SI{1}{\micro\second},
\emph{setup} de DIR de \SI{200}{\nano\second}) quedan ampliamente cubiertos por el pulso de
\SI{50}{\micro\second} generado por el firmware. Se trabaja en \textbf{paso completo}.

La corriente máxima por fase se ajusta con el potenciómetro del módulo según:
\begin{equation}
  I_{\max} = \frac{V_{\mathrm{REF}}}{8\,R_{\mathrm{CS}}}
  \qquad\Rightarrow\qquad
  V_{\mathrm{REF}} = 8\,R_{\mathrm{CS}}\,I_{\max}
  \label{eq:vref}
\end{equation}
donde $R_{\mathrm{CS}}$ es la resistencia de sensado del módulo. Los módulos utilizados tienen
$R_{\mathrm{CS}}=\SI{0,1}{\ohm}$, de modo que $I_{\max}=V_{\mathrm{REF}}/(8\,R_{\mathrm{CS}})=
V_{\mathrm{REF}}/\SI{0,8}{\ohm}$. Las tensiones de referencia ajustadas con el potenciómetro de
cada eje, y la corriente máxima por fase resultante, son:
\begin{center}
\begin{tabular}{@{}lccc@{}}
\toprule
 & \textbf{Eje X (base)} & \textbf{Eje Y (hombro)} & \textbf{Eje Z (codo)}\\
\midrule
$V_{\mathrm{REF}}$       & \SI{610}{\milli\volt} & \SI{813}{\milli\volt} & \SI{623}{\milli\volt}\\
$I_{\max}$ por fase      & \SI{0,76}{\ampere}    & \SI{1,02}{\ampere}    & \SI{0,78}{\ampere}\\
\bottomrule
\end{tabular}
\end{center}
Las referencias se ajustaron eje por eje para no superar la corriente admisible de cada motor y
limitar el calentamiento de los drivers.

\fig[0.55\textwidth]{a4988}{Módulo Pololu A4988 y su conexionado típico~\cite{pololu}.}{fig:a4988}{\texttt{Documentación/datasheets/A4988-Stepper-Motor-Driver-Pinout.jpg}}

% ------------------------------------------------------------
\subsection{Servomotor SG90 (pinza)}
El SG90~\cite{sg90} es un microservo de \SI{9}{\gram} con alimentación de
\SIrange{4,8}{6}{\volt}, recorrido aproximado de \ang{180} y torque de
\SI{1,8}{\kilo\gram\centi\meter}. Internamente contiene un motor de CC, una caja reductora y
un potenciómetro solidario al eje de salida. Un control interno compara la posición medida
con la consigna, que se codifica como el ancho del pulso de una señal PWM de
\SI{50}{\hertz} (período de \SI{20}{\milli\second}). En \tails{} se alimenta del riel de
\SI{5}{\volt} y su señal sale de PB9 (TIM4\_CH4). El ángulo se asigna linealmente entre
\SI{550}{\micro\second} (\ang{0}) y \SI{2450}{\micro\second} (\ang{180}) (ver
Sección~\ref{sec:prog-pwm}). Los ángulos de apertura y cierre son configurables por comando.

% ------------------------------------------------------------
\subsection{Finales de carrera y su acondicionamiento}
Cada eje tiene un final de carrera con conector de tres terminales (señal, \SI{3,3}{\volt},
GND) montado en un soporte impreso diseñado para el proyecto. La señal es activa en bajo:
llega a una entrada del STM32 configurada con \emph{pull-up} interno e interrupción externa por
flanco, y en la placa tiene además un \emph{pull-up} de \SI{47}{\kilo\ohm} (Figura~\ref{fig:sch-fdc}).

Cada línea gobierna también un transistor PNP \textbf{BC558}~\cite{bc558}. La base se conecta
directamente a la señal y el emisor alimenta un LED testigo desde \SI{3,3}{\volt} a través de
\SI{220}{\ohm}, con el colector a masa. Cuando el sensor se activa y lleva la línea a
\SI{0}{\volt}, la unión base-emisor conduce y el LED enciende. Con la línea en alto, el
transistor queda cortado. Así se obtiene indicación visual de cada sensor sin cargar la señal
que lee el microcontrolador. La corriente de LED resulta aproximadamente
\begin{equation}
  I_{\mathrm{LED}} \approx \frac{V_{\mathrm{CC}} - V_{F} - V_{EB}}{R}
  = \frac{\SI{3,3}{\volt} - \SI{2,0}{\volt} - \SI{0,7}{\volt}}{\SI{220}{\ohm}}
  \approx \SI{2,7}{\milli\ampere},
\end{equation}
y la corriente de base que debe absorber el sensor es $I_B \approx I_{\mathrm{LED}}/\beta
\approx \SI{15}{\micro\ampere}$ ($\beta\approx180$), despreciable.

\fig[0.35\textwidth]{sch_finales_carrera}{Etapa de acondicionamiento e indicación de un final de carrera.}{fig:sch-fdc}{\texttt{PCB/resources/sch07.png} (o \texttt{sch01.png} con los tres ejes)}

\fig[0.6\textwidth]{soportes_sensores}{Soportes de finales de carrera diseñados e impresos para el proyecto.}{fig:soportes}{render de \texttt{soporteSensor1/2.SLDPRT} o foto de los soportes montados}

% ------------------------------------------------------------
\subsection{Parada de emergencia}
Un pulsador externo conectado a PB15 (activo en bajo, \emph{pull-up}) dispara la línea EXTI15
con prioridad máxima. La acción se describe en la Sección~\ref{sec:funcionamiento}. En esta
versión la parada detiene la generación de pulsos pero mantiene el torque de retención de los
motores.

% ------------------------------------------------------------
\subsection{Display LCD 16$\times$2 con adaptador I\textsuperscript{2}C}
Display alfanumérico HD44780 de 2 líneas de 16 caracteres con un expansor de E/S PCF8574
(módulo ``LCD1602 I2C'')~\cite{lcd1602}. El controlador se maneja en modo 4 bits y cada nibble
se envía como un byte por I\textsuperscript{2}C, que incluye además las líneas RS, EN y la
retroiluminación. El firmware detecta la dirección del módulo al iniciar (0x27 ó 0x3F según
la variante del PCF8574). Si no responde, el resto del sistema sigue funcionando sin
bloquearse. Muestra el estado global, las coordenadas de los ejes y el estado de la pinza, y
también los avisos de emergencia y de conexión USB.

% ------------------------------------------------------------
\subsection{Alimentación}
La energía proviene de una fuente AT de PC, de la que se usa el riel de \SI{12}{\volt}
(\SI{11,2}{\volt} medidos en vacío). La placa tiene tres rieles (Figura~\ref{fig:sch-alim}):
\begin{itemize}
  \item \textbf{VCC (\SI{12}{\volt}):} alimentación de potencia (VMOT) de los tres A4988.
  \item \textbf{\SI{5}{\volt}:} regulador lineal \textbf{KIA7805} (TO-220)~\cite{kia7805}.
        Alimenta el servo, el LCD y la entrada del regulador de \SI{3,3}{\volt}. Tiene capacitores
        de \SI{470}{\micro\farad}\,+\,\SI{10}{\nano\farad} a la entrada y \SI{220}{\micro\farad}
        a la salida, y un diodo 1N4007 de protección inversa.
  \item \textbf{\SI{3,3}{\volt}:} regulador ajustable \textbf{LM317}~\cite{lm317}, alimentado desde
        \SI{5}{\volt}, con $R_1=\SI{330}{\ohm}$ y un preset de \SI{1}{\kilo\ohm} como $R_2$:
        \begin{equation}
          V_{O} = 1{,}25\,\mathrm{V}\left(1+\frac{R_2}{R_1}\right)
          \;\Rightarrow\;
          R_2 = R_1\!\left(\frac{V_O}{1{,}25} - 1\right) = 330\left(\frac{3{,}3}{1{,}25}-1\right)\approx\SI{541}{\ohm}.
        \end{equation}
        Alimenta la lógica de sensores y los LEDs.
\end{itemize}
Un módulo medidor de tensión y corriente intercalado a la salida de la fuente permite
monitorear el consumo (Sección~\ref{sec:resultados}). El Anexo~\ref{anx:calculos} desarrolla
la verificación térmica de los reguladores.

\fig[0.85\textwidth]{sch_alimentacion}{Etapa de alimentación: reguladores KIA7805 (\SI{5}{\volt}) y LM317 (\SI{3,3}{\volt}).}{fig:sch-alim}{\texttt{PCB/resources/sch04.png}}

% ------------------------------------------------------------
\subsection{Placa de control}
Para reemplazar el cableado sobre protoboard se diseñó en KiCad una placa que integra la
\emph{Blue Pill} (zócalo), tres zócalos para módulos A4988 con capacitores de desacople en VMOT
y puentes de selección de micropaso, conectores para los motores, los finales de carrera, el
servo, el E-STOP y el LCD, los reguladores, y LEDs de estado y testigos de tensión (VCC,
\SI{5}{\volt}, \SI{3,3}{\volt}). El ruteo usa planos de masa y pistas de mayor ancho para VCC
y VMOT. Los esquemáticos, el layout y la lista de materiales están en el Anexo~\ref{anx:pcb}.

\fig[0.7\textwidth]{pcb_3d}{Modelo 3D de la placa de control diseñada en KiCad.}{fig:pcb3d}{\texttt{PCB/resources/model.png}}

% ------------------------------------------------------------
\subsection{Estructura mecánica}
La estructura es la del MK2 Plus~\cite{thingiverse}, impresa en 3D. La base gira sobre un
rodamiento impreso de bolillas mediante una transmisión por engranajes (eje X). El hombro
(eje Y) y el codo (eje Z) se accionan con un mecanismo de paralelogramo que mantiene la
orientación del efector. A esta base se agregaron los soportes de sensores y la carcasa de la
PCB, diseñados en Solid Edge y SolidWorks.

Cada eje incorpora una reducción por engranajes rectos entre el piñón del motor y el engranaje
de salida (Tabla~\ref{tab:transmision}). Estas relaciones son las que permiten convertir los
pasos de motor en desplazamiento articular (Sección~\ref{sec:resultados}).

\begin{table}[H]
\centering\small
\caption{Relaciones de transmisión de cada eje (engranaje conducido / piñón del motor).}\label{tab:transmision}
\begin{tabular}{@{}lcccc@{}}
\toprule
\textbf{Eje} & \textbf{Piñón (motor)} & \textbf{Conducido (salida)} & \textbf{Relación $i$} & \textbf{Módulo}\\
\midrule
X (base)               & 19 & 108 & 5,68 & 1\\
Y (hombro)             & 19 & 47  & 2,47 & 1\\
Z (codo)               & 19 & 47  & 2,47 & 1\\
Pinza (\emph{gripper}) & 10 & 16  & 1,60 & 1,5\\
\bottomrule
\end{tabular}
\end{table}
